\section{Introduction}
\label{sec:introduction}

The redistricting ensemble literature has produced a rich empirical
tradition: Markov chain Monte Carlo methods generate thousands of
alternative plans, and an enacted plan's position in this distribution
characterises its extremity.
If the enacted plan's Democratic seat count falls in the 99th percentile
of the ensemble, it is described as extreme.
If it falls at the median, it is described as neutral.

This framework is descriptive, not inferential.
It does not provide p-values, confidence intervals, or power guarantees.
It does not correct for multiple comparisons when dozens of metrics are
evaluated.
It does not tell a court whether the observed extremity is statistically
significant at any specified level.

These are not abstract methodological concerns.
Courts operating under \textit{Rucho v.\ Common Cause}~\citep{rucho2019}
at the federal level have rejected partisan gerrymandering claims partly
because the plaintiffs could not identify a judicially manageable standard.
State courts, which remain the principal venue for partisan gerrymandering
litigation, are increasingly demanding formal statistical evidence rather
than intuitive percentile claims.

The expert who says ``the enacted plan is more extreme than 99\% of random
plans'' will face the question: how many plans were in the ensemble?
Were the random plans drawn from a converged Markov chain?
What is the margin of error on that 99\% estimate?
Did you correct for the fact that you tested eight different metrics?

The S-track answers these questions with a unified statistical framework.

\subsection{The State of Current Practice}

A survey of redistricting expert reports filed in state courts between
2016 and 2024~\citep{chen2013unintentional,herschlag2020quantifying}
reveals four systematic gaps:

\begin{enumerate}
\item \textbf{No convergence certification}: Expert reports rarely
  compute Gelman-Rubin $\hat{R}$ or effective sample size (ESS) for
  their ensembles.
  G.4~\citep{g4paper} shows that a typical 10,000-step NC ensemble
  has ESS $\approx 695$ for Polsby-Popper compactness---the nominal
  sample size overstates the effective sample by $14\times$.
\item \textbf{No calibrated p-values}: Percentile claims are point
  estimates with sampling uncertainty.
  An ensemble with $\ess \approx 695$ and a plan at the 0.3rd percentile
  has a 95\% confidence interval on the true percentile of approximately
  $[0.0\%, 1.2\%]$.
  The confidence interval matters: a plan at the 1.2\% boundary is
  less extreme than one at 0.0\%.
\item \textbf{No multiple testing correction}: A typical redistricting
  analysis evaluates 7--12 metrics (compactness, efficiency gap,
  mean-median, partisan bias, etc.).
  At $\alpha = 0.05$, a 7-metric analysis expects 0.35 false positives
  per state---which is small.
  Across 50 states and 3 cycles, the expected number of false positives
  grows to 52 (Section~\ref{sec:challenge3}).
\item \textbf{No power analysis}: Expert reports do not typically specify
  the probability that the ensemble would detect a genuine gerrymander
  of a given severity if one were present.
  S.3~\citep{s3paper} shows that a 10,000-step NC ensemble has 80\%
  power to detect a 10pp partisan advantage but only 45\% power for
  California.
\end{enumerate}

\subsection{Road Map}

Section~\ref{sec:challenges} analyses the four challenges in detail.
Section~\ref{sec:framework} proposes the three-component framework.
Section~\ref{sec:track} describes the S-track contributions.
Section~\ref{sec:conclusion} concludes.
